\renewcommand*{\authorOfNextLines}{Helmut Quinz}
Für \TeX -(Haupt)dokumente ist eine Dokumentklasse zu definieren.
Diese Definition steht am Beginn jedes Hauptdokuments!.
Die Dokumentklasse entscheidet unter anderem welche Makros insbesondere bezüglich der Dokumentstrukturierung zu Verfügung stehen.
Neben den von \TeX definierten Dokumentklassen lassen sich zusätzliche über Makropakete bereitgestellte Klassen mit speziellen Eigenschaften verwenden.
Besonders erwähnenswert ist dabei KOMA-Skript. 
\TeX und seine Erweiterungen sind auf den Bedarf des US-Amerikanischen Raums zugeschnitten. 
Der Entwickler Markus Kohm hat \LaTeX ~mit seinem Makro-Paket auf europäische, insbesondere deutschsprachige Verhältnisse angepasst.
Sein Paket wird aber mittlerweile aufgrund seiner Flexibilität auch gerne außerhalb Europas verwendet.\\

Für den Schulgebrauch wichtige Dokumentklassen sind: 
\begin {table} [ht]
\stretchTableRows{1.2}
\begin{tabular}{p{20mm}p{25mm}p{100mm}}
	\TeX Klasse & KOMA-Klasse 	& Anwendung\\
	article 	& scrartcl		& kleine Dokumente wie Hausaufgaben oder Protokolle\\ 
	report		& scrreprt		& Umfangreiche Dokumente wie Abschlussarbeiten oder Skripta\\
	book		& scrbook		& SEHR umfangreiche Dokumente die Buch-Stärke erreichen\\
	letter		& scrlttr2		& Briefe\\
	slides		&				& Präsentationen\\
	
\end{tabular}
\caption {Wichtige \LaTeX -Dokumentklassen}
\end {table}
\FloatBarrier
Am auffälligsten unterscheiden sich die Klassen \emph{article} \emph{report} und \emph{book} durch die Bereitstellung von Strukturierungsangaben.
Diese Strukturierungsangaben legen Dokumenthierarchien insbesondere für die Inhaltsangabe fest.
Folgende Gliederungsmerkmale stehen zur Verfügung
\begin {table} [ht]
\stretchTableRows{1.2}
\begin{tabular}{clccc}
	Level	& Ebene 					& \multicolumn{3}{c}{Verfügbar in }\\
			&							& (scr)article 	& (scr)report 	& (scr)book	\\
	-1		& \verb|\part| 				& ---			& ---			& X\\ 
	0		& \verb|\chapter|			& ---			& X				& X\\ 
	1		& \verb|\section| 			& X				& X				& X\\ 
	2		& \verb|\subsection| 		& X				& X				& X\\ 
	3		& \verb|\subsubsection| 	& X				& X				& X\\ 
	4		& \verb|\paragraph|		 	& X				& X				& X\\ 
	5		& \verb|\subparagraph|	 	& X				& X				& X\\ 
	
\end{tabular}
\caption {Gliederungsebenen in \LaTeX}
\end {table}
\FloatBarrier